% Integrated bibliography. User-supplied ref01--ref50 retain their original text.
% These supplied records have not been independently validated in this integration.
% See sources/BIBLIOGRAPHY_REVIEW.md for provenance and review flags.
\begin{thebibliography}{99}

\bibitem{ref01} Alvarez-Nodarse R: Second order difference equations for certain families of ``discrete'' polynomials, \emph{Journal of Computational and Applied Mathematics}, 99 (1998), 1--2, 37--45.

\bibitem{ref02} Alvarez-Nodarse R, Garc\'ia AG and Marcell\'in F: On the properties for modifications of classical orthogonal polynomials of discrete variables, \emph{J. Comput. and Appl. Math.}, Vol. 65, (1995), 3--18.

\bibitem{ref03} Area I, Godoy E, Ronveaux A, Zarzo A: Minimal recurrence relations for connection coefficients between classical orthogonal polynomials: Discrete case, \emph{Journal of Computational and Applied Mathematics}, 89 (1998), 2, 309--325.

\bibitem{ref04} Atakishiev M. Natig, Andr\'e Ronveaux and K. B. Wolf: Difference equation for the associated polynomials on the linear lattice, \emph{Theoretical and Mathematical Physics} 106 (1996), No. 1, pp. 76--83.

\bibitem{ref05} Atakishiev NM and Wolf KB: Fractional Fourier-Kravchuk transform, \emph{Journal of the Optical Society of America A}, 14 No. 7, 1997, 1467.

\bibitem{ref06} Atakishiev NM and Wolf KB: Approximation on a Kravchuk functions, \emph{Rev. Mex. Fis.}, 40 (1994), pp. 366--377.

\bibitem{ref07} Bose RC and Mesner DM: On linear associative algebras corresponding to association schemes of partially balanced designs, \emph{Ann. Math. Stat.}, 30 (1959) 21--38.

\bibitem{ref08} Bose RC and Shimamoto T: Classification and analysis of partially balanced incomplete block designs with two associate classes, \emph{J. Amer. Stat. Assoc.}, 47 (1952) 151--184.

\bibitem{ref09} Delsarte P: Bounds for restricted codes, by linear programming, \emph{Philips Res. Reports}, 27 (1972) 272--289.

\bibitem{ref10} Delsarte P: Four fundamental parameters of a code and their combinatorial significance, \emph{Info. \& Control}, 23 (1973) 407--438.

\bibitem{ref11} Delsarte P: An algebraic approach to the association schemes of coding theory, \emph{Philips Research Reports Supplements}, No. 10 (1973).

\bibitem{ref12} Dette H: New Bounds for Hahn and Krawtchouk Polynomials, \emph{SIAM J. Math. Anal.} 26 6 (1995), pp. 1647--1662.

\bibitem{ref13} Dunkl CF: A Krawtchouk polynomial addition theorem and wreath products of symmetric groups, \emph{Indiana Univ. Math. J.}, 25 (1976) 335--358.

\bibitem{ref14} Dunkl CF and Ramirez DF: Krawtchouk polynomials and the symmetrization of hypergraphs, \emph{SIAM J. Math. Anal.}, 5 (1974) 351--366.

\bibitem{ref15} Eagleson GK: A characterization theorem for positive definite sequences of the Krawtchouk polynomials, \emph{Australian J. Stat.}, 11 (1969) 29--38.

\bibitem{ref16} Feinsilver P and Fitzgerald R: The spectrum of symmetric Krawtchouk matrices, \emph{Lin. Alg. \& Appl.}, 235, 1996, 121--139.

\bibitem{ref17} Feinsilver P and Schott R: Krawtchouk polynomials and finite probability theory, \emph{Probability Measures on Groups X}, Plenum, 1991, pp. 129--135.

\bibitem{ref18} Foupouagnigni M, Koepf W, Ronveaux A: Fourth-order difference equation for the associated classical discrete orthogonal polynomials, \emph{Journal of Computational and Applied Mathematics}, 92 (1998), 103--108.

\bibitem{ref19} Godoy E, Ronveaux A, Zarzo A, Area I: On the limit relations between classical continuous and discrete orthogonal polynomials, \emph{Journal of Computational and Applied Mathematics}, 91 (1998), 1, 97--105.

\bibitem{ref20} Groza VA: Quantum Algebras, Kravchuk q-Polynomials, and Kravchuk-Meixner q-Functions, \emph{Ukrainian Mathematical Journal} 44 7 (1992) p. 887.

\bibitem{ref21} Ismail MEH and Simeonov P: Strong asymptotics for Krawtchouk polynomials, \emph{Journal of Computational and Applied Mathematics}, 100 2 (1998), 121--144.

\bibitem{ref22} James AT: The relationship algebra of an experimental design, \emph{Ann. Math. Stat.}, 28 (1957) 993--1002.

\bibitem{ref23} Johansson K: Discrete orthogonal polynomial ensembles and the Plancherel measure, \emph{math.CO/9906120} (43 pages).

\bibitem{ref24} Koepf W, Schmersau D: Representations of orthogonal polynomials, \emph{J. Comp. Appl. Math.}, 1998, to appear [Konrad-Zuse-Zentrum Berlin (ZIB), Preprint SC 97-06, 1997].

\bibitem{ref25} Koepf W and Schmersau D: Algorithms for classical orthogonal polynomials, Konrad-Zuse-Zentrum Berlin (ZIB), Preprint SC 96-23, 1996.

\bibitem{ref26} Koepf W: Identities for families of orthogonal polynomials and special functions, \emph{Integral Transforms and Special Functions}, Volume 5, 1997, 69--102.

\bibitem{ref27} Krasikov I and Litsyn S: On Integral Zeros of Krawtchouk Polynomials, \emph{Journal of Combinatorial Theory Ser. A}, 150 (1996), pp. 441--447.

\bibitem{ref28} Krawtchouk M: Sur une g\'en\'eralisation des polynomes d'Hermite, \emph{Comptes Rendus}, 189 (1929) 620--622.

\bibitem{ref29} Krawtchouk M: Sur la distribution des racines des polynomes orthogonaux, \emph{Comptes Rendus}, 196 (1933) 739--741.

\bibitem{ref30} Kuijlaars ABJ and Rakhmanov EA: Zero distributions for discrete orthogonal polynomials, \emph{Journal of Computational and Applied Mathematics}, 99 1--2 (1998), 255--274.

\bibitem{ref31} Levenstein VI: Krawtchouk polynomials and universal bounds for codes and design in Hamming spaces, \emph{IEEE Transactions on Information Theory}, 41 5 1995 (Sep.), pp. 1303--1321.

\bibitem{ref32} MacWilliams FJ and Sloane NJA: \emph{The Theory of Error-Correcting Codes}. The Netherlands, North Holland, 1977.

\bibitem{ref33} Pryzva GY: Kravchuk Orthogonal Polynomials, \emph{Ukrainian Mathematical Journal}, 44 7 (1992), 792.

\bibitem{ref34} Ronveaux A, Zarzo A, Area I, Godoy E: Transverse limits in the Askey tableau, \emph{Journal of Computational and Applied Mathematics}, 99 1--2 (1998), 327--335.

\bibitem{ref35} Roos B: Binomial Approximation to the Poisson Binomial Distribution: The Krawtchouk Expansion, \emph{Preprint} No. 97-13, Institut f\"ur Mathematische Stochastik, Universit\"at Hamburg.

\bibitem{ref36} Roos B: Multinomial and Krawtchouk approximations to the generalized multinomial distribution, \emph{Preprint} No. 98-6, Institut f\"ur Mathematische Stochastik, Universit\"at Hamburg.

\bibitem{ref37} Savva VA, Zelenkov VI: Multilevel systems dynamics and orthogonal polynomials of a discrete variable, \emph{Prepr.} 666, Institute of Physics, Minsk, 1992.

\bibitem{ref38} Savva VA, Zelenkov VI: Peculiarities of the dynamics of quantum systems with degenerate levels, \emph{Nonlinear Phenomena in Complex Systems}. Proc. 4th Seminar. Feb. 1996, Polatsk, Minsk, 1996.

\bibitem{ref39} Schoutens W: Levy-Sheffer and IID-Sheffer polynomials with applications to stochastic integrals, \emph{Journal of Computational and Applied Mathematics}, 99 (1998), 1--2, 365--372.

\bibitem{ref40} Smirnov YF: On factorization and algebraization of difference equations of hypergeometric type, \emph{Proceedings of the International Workshop on Orthogonal Polynomials in Mathematical Physics}, M. Alfaro et al. (Eds.), Leganes, 24--26 June, 1996. P. 153.

\bibitem{ref41} Sookoo N: Generalized Krawtchouk polynomials: new properties, \emph{Archivum Mathematicum (Brno)}, 36 (2000), 9--16.

\bibitem{ref42} Szego G: \emph{Orthogonal Polynomials}, Colloquium Publications, Vol. 23, New York, AMS, revised edition 1959, p. 35--37.

\bibitem{ref43} Tratnik MV: Multivariable Meixner, Krawtchouk, and Meixner-Pollaczek polynomials, \emph{Journal of Mathematical Physics}, 30 12, 1989, 2740.

\bibitem{ref44} Vainerman LI and Filimonova NB: Hyperspectral imagery with the application of Krawtchouk polynomials, \emph{Proceedings of the International Society for Optical Engineering}, volume 2231, Bellingham, Washington, 1994, pp. 148--155 [Software for IBM-compatible personal computers created.].

\bibitem{ref45} Vainerman LI and Filimonova NB: Signal processing and polynomials on discrete argument, \emph{Systems and Networks: Mathematical Theory and Applications}, Volume II, Math. Research Vol. 79, Akademie Verlag, Berlin 1994, pp. 889--890 [Software for IBM-compatible personal computers created].

\bibitem{ref46} Vera-Jones D: Finite bivariate distributions and semi-groups of nonnegative matrices, \emph{Q. J. Math. Oxford}, 22 2 (1971) 247--270.

\bibitem{ref47} Wielandt H: \emph{Finite Permutation Groups}, Academic Press, New York, 1964.

\bibitem{ref48} Yang B-F: Studies in finite geometries and the construction of incomplete block designs VII, \emph{Acta Math. Sinica}, 15 (1965) 812--825.

\bibitem{ref49} Yang B-F: Studies in finite geometries and the construction of incomplete block designs VIII, \emph{Acta Math. Sinica}, 15 (1965) 826--841.

\bibitem{ref50} Zeng J: Linearisation de produits de polynomes de Meixner, Krawtchouk, et Charlier, \emph{SIAM Journal on Mathematical Analysis}, Volume 21, No. 5, 1990, 1349.

% Existing verified article entry, preserved.
\bibitem{KocikFeynman}
Jerzy Kocik, \emph{Krawtchouk matrices from the Feynman path integral and
from the split quaternions}, arXiv:1604.00109 (2016),
\url{https://doi.org/10.48550/arXiv.1604.00109}.
The supplied 36-page PDF is archived as
\texttt{sources/KrawtchoukFeynman.pdf}; integration emphasizes Section~7, beginning on page 15.

\end{thebibliography}
